% -----------------------------------------------------------------------------
% This file is automatically generated.
% Do not edit manually.
% Source: notebooks/support/probability-distributions/support.Rmd
% -----------------------------------------------------------------------------



\noindent The functions \ttblue{dpois()} and
\ttblue{ppois()} are used for the calculation of
probabilities from the Poisson distribution.
They require two parameter values as arguments: the number
of errors $k$ and the value of $\mu = n\pi$.
For the \emph{Case: Number of students} remember that $n =
60$ and $\pi = 17 / 331$.


\par\addvspace{\topsep}
\begingroup
\fvset{listparameters={%
  \setlength{\topsep}{0pt}%
  \setlength{\partopsep}{0pt}%
  \setlength{\parsep}{0pt}%
  \setlength{\itemsep}{0pt}%
}}
\begin{Verbatim}[breaklines=true,formatcom=\color{ada_blue}]
poisson.pmf(1, mu=60 * 17 / 331)
poisson.pmf(3, mu=60 * 17 / 331)
\end{Verbatim}
\begin{Verbatim}[breaklines=true,formatcom=\color{ada_light_blue}]
[1] 0.1414044

[1] 0.2237979
\end{Verbatim}
\endgroup
\par\addvspace{\topsep}
